\section{回测执行与评价口径}
\subsection{周度决策流程}
每周最后一个样本内实际交易日为调仓日。程序依据交易日期分组确定该日，避免将周五是否开市写死。节假日周可能在周四或更早结束，最后一周还可能因样本终点截断。指定路径共包含 \primaryRebalances{} 次周度决策，末周和最后持有期均保留截断标记。

每次决策首先取得历史窗口，然后执行资产准入和估计，求解风险预算参考，再求解最终权重。全资产向量按固定映射还原，交易量则由目标与交易前权重的差计算。信息截止日严格早于调仓日，任何包含调仓日以后收益的窗口都不能通过验证。

\subsection{调仓日收益约定}
现有研究沿用调仓日目标权重参与当日收益的约定。权重的估计截止在前一交易日，但实际成交价格尚未建模。该约定能够明确避免以当日完整收益求解当日权重，却不等于已经证明目标可按前收盘价全部成交。模型结果应被理解为统一记账约定下的研究路径。

\subsection{交易前漂移权重}
在完整价格观察下，若上一日用于收益记账的权重为 $w_{i,t-1}$，该日资产收益为 $r_{i,t-1}$，交易前相对权重可表达为
\[
w^-_{i,t}=\frac{w_{i,t-1}(1+r_{i,t-1})}
{\sum_jw_{j,t-1}(1+r_{j,t-1})}.
\]
分子表示各资产经历价格变化后的相对市值，分母将其归一到单位预算。即使目标权重没有主动改变，不同资产收益也会导致当前权重偏离旧目标。因此，用相邻目标之差代替交易量，会遗漏维持配置所需的交易。

当前研究采用标准化权重账户，费用通过组合收益扣减。它不是逐笔现金账户，没有显式求解费用支付后的股数或整数交易单位。本文所称“无杠杆”指权重的多头单位预算性质，而非对所有真实融资、结算与现金流细节的建模声明。

\subsection{换手与成本}
调仓日换手及约定成本为
\[
T_t=\sum_i|w_{i,t}-w^-_{i,t}|,\qquad c_t=0.0003T_t.
\]
该换手统计买入和卖出金额之和。若完整卖出一项资产并将全部资金买入另一项，买卖两边都会进入 $T_t$；因此不能再将该总和直接称作单边换手。3 bps 是每单位买入或卖出金额的费率。

毛收益和净收益分别为
\[
r_t^{\rm gross}=w_t^\mathsf T r_t,\qquad
r_t^{\rm net}=r_t^{\rm gross}-c_t.
\]
非调仓日交易量为零，权重继续漂移。该成本假设仅用于可比性，不宣称是独立测量的市场冲击或全部执行费用。不同 ETF 的流动性差异、最低佣金、买卖价差及交易规模均未单独估计。

\subsection{自然周与持有期}
自然周收益由该周所有日收益复利得到。调仓后持有期则从决策日开始，到下一次决策前结束。由于主模型在周末调仓，二者的日期集合通常不同。将自然周收益直接附在本周末目标权重旁边，并将其称为该目标的后续收益，会产生解释错误。

仓库为两者分别保存字段。周度摘要含自然周起止日期、实际交易日、调仓日和持有期末日，并标明最后一期是否截断。持仓 CSV 保留全部 30 个资产行，使零权重和未准入状态也可追溯。论文图表用于展示结构，具体某一期的完整交易应回到记录核对。

\subsection{年化收益与波动}
日净值由 $V_t=\prod_{s\leq t}(1+r_s^{\rm net})$ 构造。项目现有汇总函数对截取后净值序列使用首末比值，其年化收益定义为
\[
R_{\rm ann}=\left(\frac{V_N}{V_1}\right)^{\annualizationDays/N}-1.
\]
因此首个评价日收益已进入初始净值，但在首末比值中被消去；波动和夏普亦由净值相邻变化得到。本文为保持既有五模型口径保留这一边界实现，不能将其与显式增加 $V_0=1$ 的定义混为一谈。全量日收益与周收益的复利校验仍包含首日，二者应分别核对。

设由净值恢复的日收益为 $u_t$，年化波动为
\[
\sigma_{\rm ann}=\operatorname{sd}(u_t)\sqrt{\annualizationDays}.
\]
标准差采用样本自由度。年化系数是研究约定，不意味着每个自然年实际交易日数均恰好相同。部分年度的年化收益也不同于截至该日的实际累计收益，年度表格会明确标注。

\subsection{夏普、Sortino 与回撤}
夏普比率计算如下\footnote{本文统一设定无风险利率 $r_f=0$，夏普比率以日净收益相对于零收益基准计算，Sortino 比率的最低可接受收益亦设为零。该设定为五模型比较的统计口径，不表示实际无风险资产收益为零。}：
\[
S=\frac{\operatorname{mean}(u_t)}{\operatorname{sd}(u_t)}\sqrt{\annualizationDays}.
\]
该分子为算术平均收益，不是几何年化收益。因此，不能用表中的净年化收益除以年化波动并要求得到完全相同的夏普。Sortino 用负收益的二阶矩替代标准差，正收益对应下行项为零，均值仍基于完整收益序列。

最大回撤定义为
\[
D_t=\frac{V_t}{\max_{s\leq t}V_s}-1,\qquad
D_{\max}=\min_t D_t.
\]
表格以负数表示最深下跌。回撤关注路径及先后顺序，与只观察收益分布的波动不同。同样的日收益集合若排列次序不同，最终累计收益可以相同，而最大回撤可能不同。

\subsection{尾部损失与换手汇总}
令日损失 $L_t=-r_t^{\rm net}$。项目先计算损失的 95\% 样本分位数，再对不低于该分位数的损失取平均，作为报告的历史 CVaR。有限样本中的分位数插值及边界并列值，会使入选观察数量不必恰好等于样本的 5\%。正文的 CVaR 均沿用这一实现。

月均换手为全部调仓换手之和除以评价区间覆盖的自然月数，而非每次交易的平均值。年化收益成本拖累定义为同一路径毛、净年化收益之差，包含复利影响，不等于简单将费率乘以自然年换手。分别定义这些量，可以防止在图表解释中多扣一次费用或混用年化方式。

\subsection{验证链条}
独立核验读取保存的权重、收益、换手和费用，重新计算净值及汇总指标。预算、非负性、准入和信息截止日期逐项检查。年度参数表还需满足验证期末早于参数生效日，避免将当年收益用于当年参数选择。

未来扰动和晚上市资产测试针对输入时序；周度导出测试针对收益、成本与持有期衔接。它们覆盖具体实现行为，不证明统计显著性，也不证明模型选择未使用历史结果。将工程正确性与研究有效性分别报告，才能准确理解验证结论。
